\chapter{Roadmap and Open Gaps}
\label{ch:roadmap}

This closing chapter gathers the open threads left throughout both volumes into one place,
organized the same way CNA's own README states its roadmap: incremental API coverage first,
then the specific named gaps, then cross-platform strengthening.

\section{Continued API coverage, class by class}

The largest item on CNA's own roadmap is also its least dramatic: continuing to expand XNA
API coverage and behavior parity incrementally, class by class, rather than in one large
push. Chapter~\ref{ch:what-is-cna} already gave the headline numbers --- 227 of 245 public
FNA types present (92.7\%), with \texttt{.Content} (four of twelve types, by design, per
Chapter~\ref{ch:content-and-assets}) and \texttt{.Media}'s fourteen previously-shell types
(now largely closed, per Chapter~\ref{ch:media-system}'s own correction to that figure) as the
two namespaces furthest from full presence. This kind of gap closes the way most of this
book's own bug histories closed: one class, one audit, one fix at a time, tracked through the
\texttt{plan\_*.md} / \texttt{AUDIT.md} mechanism from Chapter~\ref{ch:project-practice}.

\section{The single biggest named gap: compiled .fx bytecode}

Chapter~\ref{ch:shader-fx-gap} covered this in full, and it remains, in the
project's own words, the single biggest real gap in CNA: no support for loading a compiled
\texttt{.fx} shader bytecode blob the way a real XNA content pipeline would produce one, with
a real, measured cost of 23 of the 86 official XNA samples in \texttt{cna-samples}
(Chapter~\ref{ch:samples-and-examples}) failing to run as a direct result. The
\texttt{ShaderEffect} runtime-HLSL-compilation path on the three Direct3D backends
(Chapter~\ref{ch:direct3d-backends}) is a real, working alternative for new or adapted
shaders, but is not bytecode compatibility, and closing this gap for real would mean building
something comparable in scope to a second MojoShader-style bytecode cross-compiler.

\section{The .xnb content pipeline: a design choice under active reconsideration}

Chapter~\ref{ch:content-and-assets} covered CNA's raw-asset-plus-JSON-descriptor content
model as a deliberate design choice, not an oversight --- and the roadmap frames a real,
binary-compatible \texttt{.xnb} reader as something the project is actively considering
rather than something ruled out permanently. Chapter~\ref{ch:models-meshes}'s account of the
real, already-shipped \texttt{.xnb} \cnaclass{Model} reader is a concrete sign of where this
is already heading incrementally, class by class, rather than as one large content-pipeline
rewrite.

\section{Named architecture-decision gaps}

Four specific gaps recur across this book's backend chapters, each requiring a project-owner
decision rather than more engineering effort on its own:

\begin{itemize}
  \item \texttt{SDL\_RENDERER}'s \cnaclass{TextureAddressMode::Wrap} / \texttt{Mirror} via
        \cnaclass{SpriteBatch} (Chapter~\ref{ch:sdl-renderer}) --- blocked on choosing between
        three concrete implementation strategies, each with a real tradeoff.
  \item \texttt{SDL\_RENDERER}'s \cnaclass{Texture3D} / \cnaclass{TextureCube} construction
        (Chapter~\ref{ch:sdl-renderer}) --- blocked because the safer fix (throwing at
        construction rather than silently succeeding) has a confirmed 94-test blast radius.
  \item \texttt{EASYGL}'s real, non-\texttt{Color} \cnaclass{SurfaceFormat} GPU forwarding
        (Chapter~\ref{ch:easygl-backend}) --- blocked because it conflicts with an
        already-shipped \texttt{Texture::ValidateFormat} contract the project has already
        committed tests against.
  \item \cnaclass{Texture3D} / \cnaclass{TextureCube}'s sampler-bind architecture
        (Chapter~\ref{ch:textures-rendertargets}) --- the structural deviation from FNA
        where neither type participates in the generic texture-slot binding mechanism at
        all, so only stock effects with a dedicated field (\cnaclass{EnvironmentMapEffect}'s
        cube map) can use one.
\end{itemize}

\section{The last item, compared across all seven backends}

The \cnaclass{Texture3D} / \cnaclass{TextureCube} sampler-bind gap above is worth seeing at its
full, real width rather than as a single bullet, since the project's own internal feature
matrix
(\texttt{docs/\allowbreak{}graphics-\allowbreak{}backend-\allowbreak{}feature-\allowbreak{}matrix.md})
tracks it across every backend this
book has covered, not just the two it happens to name most often:

\begin{center}
\small
\begin{tabular}{lp{9.3cm}}
\toprule
\textbf{Backend} & \textbf{Texture3D/TextureCube sampled in a shader} \\
\midrule
\texttt{EASYGL} & Blocked --- neither type inherits \cnaclass{Texture} at all; this is the
  architectural root cause every other backend's own status traces back to. \\
\texttt{VULKAN} & Same architectural block. \\
\texttt{BGFX} & Same architectural block. \\
\texttt{SDL\_RENDERER} & Not applicable --- a 2D-only backend with no 3D-only resource types
  to begin with. \\
\texttt{D3D9} & Partial --- \cnaclass{TextureCube} is sampled and oracle-verified byte-exact
  via \cnaclass{EnvironmentMapEffect}; \cnaclass{Texture3D} has no consuming shader variant
  tested at all. \\
\texttt{D3D11} & Partial, the same shape --- \cnaclass{TextureCube} pixel-verified via a real
  environment-map scene; \cnaclass{Texture3D} untested in any shader. \\
\texttt{D3D12} & Partial, the same shape again, though \cnaclass{Texture3D} at least has a
  real, working backend implementation behind it (Chapter~\ref{ch:direct3d-backends}) --- it
  is specifically the shader-sampling path that remains unverified, not the resource itself. \\
\bottomrule
\end{tabular}
\end{center}

The pattern across all seven rows is itself the finding: this is not a gap unique to one
backend's own maturity, but a single architectural decision (Chapter~\ref{ch:textures-rendertargets}'s
account of why neither type participates in the generic texture-slot binding mechanism)
propagating an identical shape of limitation into every backend that implements it, with
each Direct3D backend's own \cnaclass{TextureCube} support proving the underlying resource and
GPU sampling both already work correctly --- the missing piece is specifically the generic
binding path, not GPU capability.

\section{Strengthening cross-platform execution and validation}

The roadmap's final, broadest item is exactly what this Part's cross-platform chapters
covered in depth: deepening validation on Windows (moving from Wine-verified to genuinely
Windows-hardware-verified, an open item on every one of the three Direct3D backends), the
Web (moving \texttt{CANVAS} from code-reviewed to actually pixel-verified in a real browser,
Chapter~\ref{ch:canvas-ascii-backends}), and Android (broadening beyond the single real
x86\_64 emulator configuration this project currently verifies against).

This item is not purely aspirational --- it already has a real, concrete precedent worth citing
as evidence the strategy works, not just as a plan. The Android sensor bridge's own
bounded-timeout bug (Chapter~\ref{ch:android-ndk}) was invisible to careful manual review and
to the ordinary test suite alike; it was found only once a dedicated \texttt{devices-tsan} run
was specifically constructed against real, concurrent \texttt{Start()} / \texttt{Stop()} traffic
under load. The same is true of \texttt{sharp-runtime}'s own \cnaclass{Channel} deadlock
(Chapter~\ref{ch:sharp-runtime-namespaces}) and the D3D9 backend's real swap-chain format restriction
(Chapter~\ref{ch:windows-wine}), each surfaced only once the corresponding platform's own
verification loop was deepened past what a prior pass had considered sufficient. ``Strengthen
cross-platform validation'' is not a vague aspiration in this project's own history --- it is a
strategy already shown, independently, three separate times, to find real bugs no earlier pass
had caught, which is the strongest argument available for continuing to invest in it rather
than treating current coverage as good enough.

\section{A gap in this book itself, closed: \texttt{SDL\_GPU} now has its own chapter}

An earlier pass through this chapter named a gap in this book itself, not in CNA's own code:
\texttt{SDL\_GPU}, a real, extensively-verified backend, had never received a dedicated Part~IV
chapter despite maturing alongside every other backend this book covers. That gap is closed as
of this pass --- Chapter~\ref{ch:sdlgpu-backend} now gives it the same full treatment every
other backend receives (the Y-flip bug a real screenshot caught, the storage-buffer workaround
for a real undocumented push-uniform cap, the runtime \texttt{libshaderc} compile, and the
render-target use-after-free closed before it ever shipped). This is left here, briefly, as a
tracking pointer rather than deleted outright, the same convention
Chapter~\ref{ch:project-practice} already describes for a superseded task row in this
ecosystem's own \texttt{plan\_*.md} files --- a reader who remembers this chapter naming the gap
should not be left looking for it here once it no longer exists, but pointed forward to where
the real content actually landed.

\section{A closing note on how to use this chapter}

Every gap named in this chapter is named because a specific, dated source said so ---
\texttt{NEXT.md}, a \texttt{plan\_*.md} file, or a backend's own documentation, per the
practice conventions from Chapter~\ref{ch:project-practice}. As with every dated claim this
book has made, the right way to use this chapter is not to trust it indefinitely, but to use
it as a map: go to the source it cites, check whether the gap is still open, and update your
own understanding from there. That is the same discipline this project asks of itself, and
the same discipline this book has tried to model throughout both volumes.
